% Chapter 1: the Elenchus method, explained simply.
\xdef\labelcourant{ch:elenchus}
\chapter{The Elenchus method, explained simply}

If I had to explain to someone who is not a scientist what I did this year, I would not start with the theorems. I
would start with a question anyone can ask: \emph{how can you tell whether a machine is telling the truth?} AI agents
write quickly and well. They also write false things, with the same confidence. This chapter tells the story of the
method we ended up adopting so that we no longer take them at their word. There is nothing magical about it. It comes
down to a few rules that a cook, a judge or a bridge inspector already knows.

\section{A Greek name for a simple idea}

\emph{Elenchos}, in ancient Greek, means examination by questioning, the kind Socrates practised. Socrates was not
trying to prove that he was right. He questioned the person in front of him until he found the place where their
claims contradicted each other. What withstood the examination deserved to be believed, at least for a while; the
rest had to be abandoned.

We took up the idea and turned it into engineering. Instead of a philosopher asking questions, there are \emph{rules}
that every result must get through before it is published, and \emph{programs} that ask some of the questions for us.
The method was written down in September 2026, in a paper that describes how a small AI-assisted laboratory works,
from choosing the questions all the way to publication.\reference{``A laboratory method for AI-assisted science'',
18 September 2026. \doi{10.5281/zenodo.22829767}} It was then applied, more and more strictly, to every contribution
in this book.

Here are the rules, one by one.

\begin{encadre}[The seven rules of Elenchus, on one page]
\begin{enumerate}[leftmargin=1.6em]
\item \textbf{Write down the claim before the test}, and say in advance what would bring it down.
\item \textbf{Give every claim a level}: A, L, C or X, depending on who checked it.
\item \textbf{Run a test that must fail and a test that must succeed}, to check that the test works.
\item \textbf{Keep the one who makes separate from the one who checks.}
\item \textbf{Publish the negative results}, with the same care as the others.
\item \textbf{Hand the logic to an impartial judge}, the Lean kernel, and know how it can be fooled.
\item \textbf{Keep a ledger} of every claim, with its level and its evidence, and let it say ``no'' when something is
missing.
\end{enumerate}
\end{encadre}

\section{Rule 1: write down the claim before the test}

A dishonest archer shoots his arrow first and then paints the target around it. He always hits the bull's-eye. In
science the temptation is the same: you do a calculation, you look at the result, and then you decide afterwards what
you were looking for. The remedy has long been known in medicine, where it is called \emph{preregistration}: before
running the experiment, you write down, dated and archived, what you predict, and what, if you observe it, will prove
you were wrong.

In this project, every important hypothesis was preregistered in a file, before the calculation. The file states the
claim, the probability the author gave it of surviving the test, and the ``stopping rule'': what would kill it. For
example, the conjecture we shall call QU-1 (it comes back later in this chapter) was recorded with an estimated
chance of survival of 0.6, that is, a little better than even.

Preregistering does not guarantee that you are not wrong. It guarantees that you cannot \emph{tell yourself} that you
were not wrong. And it forces an extra kind of honesty: when the real work departs from what was announced, you have
to say so. Several papers in this book have a section called ``deviations from the preregistration''. One of them, in
Chapter~\ref{ch:c07}, reports that one of its own preregistered sub-hypotheses was simply false.

\section{Rule 2: every claim carries a level}

The second rule is the most visible one in this book: every claim carries a label that says \emph{who checked it}.
There are four levels, and each has its own picture.

\begin{figure}[htbp]
\centering
\begin{tikzpicture}[
  niv/.style={rectangle,rounded corners=3pt,minimum width=5.3cm,minimum height=1.55cm,text width=4.9cm,
    align=left,inner sep=6pt,font=\small},
  lettre/.style={circle,minimum size=0.95cm,font=\sffamily\bfseries\large,text=white,inner sep=0pt}]
  \node[niv,fill=okvert!14!ivoire,draw=okvert,line width=0.8pt] (A) at (0,0)
    {\lettreniv{okvert}{A}\textbf{``The judge has looked.''}\\ A program, the Lean kernel, checked every logical step.};
  \node[niv,fill=okbleu!12!ivoire,draw=okbleu,line width=0.8pt] (L) at (5.75,0)
    {\lettreniv{okbleu}{L}\textbf{``It is written in a serious book.''}\\ A published result, read and cited precisely.};
  \node[niv,fill=okorange!14!ivoire,draw=okorange,line width=0.8pt] (C) at (0,-1.95)
    {\lettreniv{okorange}{C}\textbf{``A line of reasoning.''}\\ Carried out here by AI agents; nobody has checked it by machine.};
  \node[niv,fill=okpourpre!14!ivoire,draw=okpourpre,line width=0.8pt] (X) at (5.75,-1.95)
    {\lettreniv{okpourpre}{X}\textbf{``A pencil sketch.''}\\ A calculation with decimal numbers: it points the way, it proves nothing.};
\end{tikzpicture}
\caption{The four confidence levels, with the everyday words used throughout the book.}
\label{fig:niveaux}
\end{figure}

\paragraph{\niveau{A} ``the judge has looked''} A program, the Lean \emph{kernel}, has checked every one of the logical
steps of the proof. It is not the computer that ``thinks''; it checks that no step is skipped. To earn this level,
three checks must pass: the proof uses only three basic rules of logic, always the same ones (their technical names
are \texttt{propext}, \texttt{Classical.choice} and \texttt{Quot.sound}); a second program has reread the proof from
end to end, which is called the \emph{replay}; and a search for the shortcuts that make it possible to get round the
judge found nothing. When the replay is missing, the book writes \niveau{As}.

\paragraph{\niveau{L} ``it is written in a serious book''} The result has been published in the scientific
literature, and we have read it and cited it with page and line. We have not redone the proof, but others reviewed it
before us.

\paragraph{\niveau{C} ``a line of reasoning, carried out here by AI agents, that nobody has checked by machine''}
This is a written argument, often correct, that nobody has submitted to the judge. In the sciences, such an argument
is usually a researcher's; in this book, it was almost always carried out by one AI agent and reread by another. That
is one more reason to be careful.

\paragraph{\niveau{X} ``a pencil sketch''} This is a calculation done by computer with decimal numbers, like the ones
on a pocket calculator. These numbers are rounded at every operation; they give a very good idea, and sometimes a
wrong one. A simulation, however huge, remains a sketch.

\medskip
A complete claim often mixes several levels. In that case, \emph{it is only as strong as its weakest link}. An example
at the end of the chapter will show this.

\begin{analogie}
Think of the labels on a jar of jam. ``Home-made'' does not say the same thing as ``inspected by the food safety
office'', nor as ``grandmother's recipe, never tested''. No label is shameful; what would be dishonest is to stick
``inspected'' on a jar that nobody has opened. The picture stops there: an inspected jar can still be bad, and a
level~A theorem is only worth as much as its statement, and the judge itself can, very rarely, be wrong; it is the
question of the statement that stays open, as we shall see.
\end{analogie}

\section{Rule 3: a test that must fail, and a test that must succeed}

A smoke detector that never goes off is very reassuring, right up to the day of the fire. That is why you press its
test button: you check that it \emph{can} go off. Scientists call this \emph{controls}. A positive control is a case
where you know the answer must be ``yes''; a negative control, a case where you know it must be ``no''. If the test
answers ``yes'' to both, it is worthless.

In the autumn work, no measurement is reported without its two controls. A few examples that will come back in the
chapters:
\begin{itemize}
\item To check that the Lean judge really refuses an unfinished proof, we deliberately submit a proof \emph{with a
hole} in it, marked ``to be finished'' (in Lean, this word is written \texttt{sorry}). The control succeeds when the
hole is detected.
\item To check a calculation certificate, we also give it a case where it \emph{must} fail. In the example at the end
of the chapter, the same method, applied to another parameter where the expected result is different, ``fails, as it
should''.
\item To check an auditor, we hide mistakes in the work we hand over. In Chapter~\ref{ch:c07}, the auditor planted
mistakes in a copy of the files, and they were caught. In Chapter~\ref{ch:c13}, a deliberately damaged copy of a large
certificate was rejected.
\end{itemize}
Even the cover of this book has its controls: the program that draws it checks that every cell has six sides and an
area of~1, and then checks that this test fails on a slightly deformed lattice.

\begin{analogie}[A picture: the smoke detector]
A test that has never failed has never been tested. Before believing an ``all is well'', you check that the device
knows how to say ``watch out''. The picture has a limit: a detector that sounds during the test can still miss a fire
of another kind. A control shows that the test works on \emph{this} case, not on all of them.
\end{analogie}

\section{Rule 4: the one who makes is not the one who checks}

The cook who tastes his own sauce always finds it to his liking: he knows what he meant to put in it. That is why you
ask someone else to taste it, someone who does not know the recipe and has no reason to please the cook.

In the work of the end of the year, the AI agent that produces a result is not the one that checks it. A second agent, launched
separately, receives the files, recompiles them, redoes the calculations with its own programs and gives a written
verdict in an audit file. The verdicts have precise names: ``accepted with corrections'', ``confirmed with
corrections''. Then the author applies the corrections, and a second file, the revision file, keeps a record of them.
For example, the audit of the distances between galaxies (Chapter~\ref{ch:c01}) found thirteen defects of attribution
and scope, corrected before publication; the audit of one of the last papers of the year, on screened forces, led to twelve
revisions.\reference{\doi{10.5281/zenodo.23663551}}

This separation has an obvious limit, and we say so: the taster is also an AI, often from the same maker. It does not
know the recipe, but it may have the same taste buds. That is one reason why an audit by an AI is not a human review
(we come back to this below).

\begin{analogie}[A picture: the recipe tasted by someone else]
You do not ask the cook whether the dish is good; you have it tasted by someone who did not cook it, and you write
down what they say. The picture stops where science begins: here the taster does not just give an opinion, it makes
the recipe again on its own and compares the two dishes, gram by gram.
\end{analogie}

\section{Rule 5: we publish the negative results}

A negative result is an idea that does not work: a formula that the data contradict, a refuted conjecture, a method
that does not extend. It is very tempting to put them away in a drawer. But a drawer full of hidden failures distorts
the picture everyone has of a subject: you only see what succeeded.

This book does the opposite. Three of the twenty contributions are ``lessons'' in the strict sense, and many others
contain a refutation. In Chapter~\ref{ch:c05}, measurements on superfluid helium contradict a formula that the
project itself had proposed a few weeks earlier. In Chapters~\ref{ch:c12} and~\ref{ch:c11}, a beautiful idea (``the
triangular lattice is also the best in the quantum world'') is disproved, first by calculation, then with a
certificate. In Chapter~\ref{ch:c20}, a claim from the batch of 19 October is withdrawn because the measurement
contradicts it.

A negative result done well is not a failure: it is a map that says ``do not go this way''. It often teaches more than
the success we were hoping for.

\section{Rule 6: an impartial judge, and how it can be fooled}

\subsection{The judge}

Lean is a language in which mathematical proofs are written in so much detail that a computer can check them. At the
centre of Lean is a small program, the \emph{kernel}, whose only job is to check each step: does this one really
follow from the previous one, according to the accepted rules? The kernel cannot invent a proof. It can only say
``yes, every step is correct'' or ``no, something is wrong here''.

It is a marker that does not tire, does not flatter and is rarely wrong. It does not tire: it rereads the hundredth
page with the same attention as the first; the audit in Chapter~\ref{ch:c15} had it reread 88,400 declarations. It does
not flatter: it does not know who wrote the proof, a famous AI or a beginner, and it does not care. It is rarely wrong:
its code is small and public, designed so that it can be reread; but rarely does not mean never.

\begin{analogie}[A picture: the trial]
Imagine a trial. The AI agent that produced the result is the defence: it argues that it is true. The agent that
audits it is the prosecution: it looks for the flaw. The Lean kernel is the judge, and this judge has one peculiarity:
it reads only the written evidence and the law, never the speeches. Fine words do not impress it. But the judge does
not decide everything. It does not say whether the trial was about the right question: it is up to the human, at the
end, to decide what the verdict means and whether it deserves to be published. The picture stops there: a real judge
weighs testimony, whereas the kernel weighs nothing, it checks.
\end{analogie}

\subsection{Four ways to fool the judge, and how we notice}

The judge is impartial, but it is not safe from trickery. The year taught us four tricks, and each one gave us a rule.

\paragraph{The homework handed in with a page ``to be finished''} In Lean, you can write \texttt{sorry} in place of
part of a proof: it is a way of saying ``I will come back to this later''. The program then agrees to compile the file,
and it \emph{finishes without an error}. During the summer, one paper claimed that its Lean checks had succeeded
because the program had finished without an error. But a \texttt{sorry} also finishes without an error. The remedy:
never trust a mere ``finished without an error'', but ask the judge for the list of what each theorem rests on. If a
hole is hiding in it, it shows up under the name \texttt{sorryAx}.

\paragraph{The rule slipped in on the quiet} In Lean, you can also declare an \emph{axiom}: a rule that is accepted
without proof. If you accept a false rule, you can prove anything, including $1=2$. One of the pieces of work in this
book showed this on purpose (Chapter~\ref{ch:c18}): a file ``proving'' $1=2$ thanks to a hidden axiom compiles without
the slightest error, and only the audit of the axioms unmasks it. Hence the rule: the list of axioms of every theorem
must be limited to the three basic rules mentioned above.

\paragraph{The document slipped into the file without going before the judge} This is the subtlest trick, and it was
found by our own audits. Lean allows, through an option reserved for developers, a declaration to be added
\emph{without having it checked} by the kernel. During the audit of OpenAI's proofs about the equations of fluids
(Chapter~\ref{ch:c15}), the auditor deliberately made such a document, which ``proves'' an absurdity, and observed a
negative result that it put at the top of its report: \emph{the list of axioms does not detect this
forgery}.\reference{Audit A2, 26 October 2026, \nolinkurl{AXIOMS.md} and \nolinkurl{SOUNDNESS_SCAN.md}.} The list is
clean; the judge has simply never seen the document.

The counter-measure is the \emph{replay}: a second program, \texttt{leanchecker}, rereads from the beginning every
declaration that has been recorded and submits it to the kernel again. On 26 October, at the same time as the replay of
our own modules, we repeated the experiment on a control forgery. The list of axioms answered that the document
``depends on no axioms'': the usual check would have let it through. The replay, on the other hand, stopped with an
error, with the message ``the type of the declaration does not match''.\reference{Replay of 26 October 2026,
\nolinkurl{results/kernel_replay_2026-10-26/}, control \texttt{KRForged}.} The replay was run on 93 of our proof
modules, and all of them passed; a search for the four known bypass options in 871 files found nothing.

\begin{analogie}[A picture: the bridge inspection]
A bridge carries a plate: ``inspected and compliant''. The plate is the list of axioms: it says what the file
contains. But a file can contain a page that the inspector never signed. The replay is the inspector coming back and
going round the bridge again, bolt by bolt, without looking at the plate. The picture has a limit: a bridge wears out,
a proof does not. Once the replay has passed, there is no need to do it again, unless the files or the tools change.
\end{analogie}

\paragraph{The right proof of the wrong question} The last trick is not really a trick, and it is the most important.
The judge checks that the proof proves the written statement. It does not check that the written statement says what
we \emph{tell people} it says. If I write in Lean ``every square is a rectangle'' but get the definition wrong, the
judge will happily validate a theorem of no interest. Only a reader who compares the formal statement with the
question that was asked can catch this kind of mistake. This step, which we call the \emph{audit of the statement},
has not yet been done by any human for the results in this book. The ledger, the next rule, says so every time.

\section{Rule 7: the ledger of claims}

Every paper from the end of the year comes with a \emph{ledger}: a file that contains one line per claim. Each line
carries an identifier, the level, the statement in words, a digital fingerprint of the proof (a kind of barcode that
changes if the smallest file is modified), and a ``human audit'' box. Here, translated, is a real line from the
ledger of the audit of the proofs about fluids (Chapter~\ref{ch:c15}):

\begin{encadre}[One line of the ledger]
\small
\textsf{\textbf{P2NS-A-0001}} \quad \niveau{A} \quad \textsf{human audit: none}\par\smallskip
\emph{Statement:} the Lean theorem \texttt{navier\_stokes\_breakdown\_R3} states that, for every positive viscosity,
there exist an initial velocity and a force, smooth and decaying very fast, for which no smooth solution of bounded
energy of the Navier--Stokes equations exists. Our own audit of the axioms gives exactly the three basic rules; the
module was rebuilt from the sources, and its 53,752 declarations were replayed by the kernel. \emph{A statement about
the Lean term only.}\par\smallskip
\emph{Proof:} fingerprint \texttt{sha256:0c86f9d7\dots}
\end{encadre}

A program, the \emph{ledger gate}, reads all the lines and refuses to say ``all is well'' as long as something is
missing. For this paper, its answer is: four times the signal \texttt{LEDGER\_UNAUDITED\_TIER\_A}, ``the kernel has
checked the proof, and no person has yet certified that the statement means what its gloss says'', then exit
code~1, which means ``no''.

That is intentional. Almost all the ledgers in this book end this way, because the human audit of the statements has
not been done. The author chose to publish anyway, and to write so in every paper. This is already the case for the
paper in Chapter~\ref{ch:c03}, whose ledger lists twenty-two claims. Publishing with the red flag raised, and showing
it, is more honest than lowering the flag.

\section{What an AI audit is not}

It has to be said clearly, because the year caught us out on this point: \textbf{an audit by an AI agent is not peer
review}. Peer review, in science, is the examination of a piece of work by human researchers in the field, who stake
their judgement and their reputation on it. Nothing in this book has received that examination.

On 19 and 20 October, papers appeared with the label ``peer reviewed in five rounds''. On 21 October, one of them
published a second edition that withdraws this label: the attached review file ``was produced by a script that calls
no reviewer''.\reference{Withdrawal of 21 October 2026: \doi{10.5281/zenodo.23525655}} As early as 19 October, another
record had withdrawn a ``review report'' written by a program, with imaginary reviewers.\reference{\doi{10.5281/zenodo.23468458}}
A month earlier, the audit of the workshop itself had found a ``peer review'' script that never called a model.

Hence two rules, which we now apply to the letter. A review file must show that a reviewer was really called: who,
when, and with what answer. And we write ``audited by an AI agent'', never ``peer reviewed'', as long as no human in
the field has done it.

\section{An example: an idea that goes up and down the levels}

Here, to finish, is the story of a single claim, from rough draft to ledger. It is told in detail in
Chapters~\ref{ch:c12} and~\ref{ch:c11}; here we follow only its levels.

\paragraph{The idea} We know (it is the thread of the second part) that the triangular lattice minimises the energy
of points that repel one another when they are \emph{at rest}. But in the quantum world, points are never completely
still: even at absolute zero, they vibrate a little. This minimal vibration has an energy, the \emph{zero-point
energy}. The natural idea, christened QU-1, was: ``the triangular lattice also minimises the zero-point energy''. At
this stage it is a conjecture: not even an argument, an intuition. It is preregistered, with an estimated chance of
survival of~0.6.

\paragraph{The pencil sketch \niveau{X}} Before trying to prove it, we test it by calculation. Two formulas, programmed in four ways
and then a fifth time by the auditor, compute the zero-point energy of many lattices. Result: for eight of the twelve forces tested, another
lattice, a stable one, has a lower zero-point energy than the triangular lattice. QU-1 is refuted, but only ``in
pencil'': decimal numbers, however well they agree, prove nothing.\reference{Numerical refutation, 26 October 2026:
\doi{10.5281/zenodo.23653833}}

\paragraph{The certificate \niveau{A} for the logic, \niveau{C} for the whole} For one of the forces, the Gaussian bell,
we want a proof. We deform the triangular lattice very slightly, by a small \emph{shear}, the way you push the top of
a stack of cards. A rigorous calculation by enclosures, which guarantees at every step that the true value lies in an
interval (this is called ball arithmetic), shows that the zero-point energy \emph{decreases} for every small enough
shear: so the triangular lattice is not even a local minimum. The calculation split the problem into 206,468 cells,
and none of them failed. The Lean kernel checks the logic of the argument and the exact arithmetic of the
certificate.\reference{Certificate, 26 October 2026: \doi{10.5281/zenodo.23655392}}

But the enclosures themselves were computed by a program \emph{outside} the kernel, and some of the steps that link
the pieces together are reasoning written by AI agents, which neither a machine nor a human has checked (level~C). So the paper writes, honestly: ``the main result is level~C as a whole;
only its logic and its arithmetic are level~A''. That is the weakest-link rule.

\paragraph{The controls} The same method, applied to a narrower bell for which the ``pencil'' calculations do not
predict a decrease, does not manage to certify a decrease, ``as it should''. Proofs with deliberate holes, a planted
axiom and a falsified certificate were submitted to the judge, which refused them all.

\paragraph{The audit and the ledger} A distinct agent redid the calculations; a later paper extended the result to
screened forces, with an audit and twelve revisions.\reference{Screened forces, 26 October 2026:
\doi{10.5281/zenodo.23663551}} And the ledger gate still answers ``no'': no person has yet certified that the Lean
statements say what their glosses say.

\begin{encadre}[What the example shows]
Within a few weeks, the same idea was an intuition, a preregistered hypothesis, a refutation in pencil, and then a
mixed certificate in which every piece carries its level. At no point did we write ``proved'' for something that was
not. That is Elenchus: not always being right, but knowing at every moment \emph{how far} you are right, and saying
so.
\end{encadre}

\section{Even the biggest: checking another AI's proof}

The rules do not apply only to our own agents. They apply, just as strictly, to the work of the big laboratories. In
2026, OpenAI published proofs produced by an AI and checked by computer. It would have been tempting to take
them as they were: they came from a renowned team, and the Lean judge had accepted them. Even so, we examined them as
we examine our own, before building on them.

\paragraph{A claim about electric charges} The publication also claims to prove the Sandier--Serfaty conjecture, about
the most economical arrangement of charges in a neutralising background, by a reduction and a very large calculation
certificate. An agent ran this certificate again: it reproduced in about eight minutes. It checked that a deliberately
damaged copy is rejected (that is the negative control), and it read the reduction step by step without finding a
gap. Chapter~\ref{ch:c13} tells this story, and its title says it: \emph{no proof is claimed}. Neither proved, nor
verified, nor refuted by us; reproduced and audited, which is different. Along the way, the audit found and repaired a
gap in \emph{our} certificate tool. Then four of the twelve stages of the certificate were redone inside the kernel
(Chapter~\ref{ch:c14}); the other eight are still waiting.

\paragraph{Theorems about fluids} OpenAI also published, in Lean, ``blow-up'' theorems for the Navier--Stokes and Euler
equations, which describe how fluids flow. An agent rebuilt them independently: 2,440 modules recompiled, 88,400
declarations replayed by the kernel, no loophole found, and the statements compared with the official texts of the
problems (Chapter~\ref{ch:c15}). This work does not claim that a Millennium Prize Problem has been solved, and the
credit for the proofs belongs to their authors. It only says, in a way anyone can check, what was checked and what
was not.

\begin{analogie}[A picture: the roadworthiness test]
A new car from a big manufacturer still goes through its roadworthiness test. This is not a suspicion against the
manufacturer; it is the rule, the same for everyone. The picture stops there: the inspector of a car knows the list
of points to check, whereas for a new proof, part of the work consists in inventing that list.
\end{analogie}

\section{The levels of the twenty, at a glance}

The table at the end of the book gives, for each contribution, its rigour score in the selection grid, from 0 to 3.
Sixteen score 3: a proof checked by the judge and replayed, or a calculation certificate rechecked by a different
program. One scores 2 (checked by the judge, without replay), two score 1, and one scores 0: the calculation
instrument rusty-SUNDIALS, all of whose numbers are decimal numbers. That is intentional: an instrument does not have
to be a theorem.

This count has to be read with care. The score is that of the \emph{best} piece of each work, not of the work as a
whole; nor is it quite an Elenchus level, since a certificate rechecked without a replay by the judge counts as much
as a replay. The level badges in each chapter are what count. A level~A contribution almost always also contains
simulations (X), citations (L) and arguments (C), and its main conclusion is only as strong as its weakest link, as
in the example of QU-1. That is why every chapter has a section ``How we made sure it is not wrong'', which goes
through the levels piece by piece, and a section ``What this does not say''.

\section{The rules applied to this book}

This book tries to follow its own rules. Its sources are fixed in advance: a number that does not appear in the file
of a piece of work is not used here. The twenty contributions were chosen with a grid that was written and frozen
before a single score was given, and applied in the same way to every candidate (this is the equivalent of
preregistration). Every chapter follows the same plan, whatever the field, so that no result is presented better
than another simply because it is more appealing. The mistakes and the withdrawals have their own chapter at the end.
And this book, like its results, has not been reviewed by an outside human scientist: so it, too, carries a raised
flag.

\section{What the method does not do}

Elenchus is not a guarantee. It does not replace human review; it does not say whether a question is worth asking; it
does not prevent mistakes, it makes mistakes easier to find. And it does not apply itself: the relapse of 19 October
came one month \emph{after} the method had been written down. A rule on paper protects against nothing. What protects
you is running it, every single time, including when you are in a hurry and pleased with yourself.

The twenty chapters that follow are its application. Read the level before the sentence.
